% report.tex — RB GenAI evaluation, 2-page two-column report (plan §8).
% Compile: pdflatex report && bibtex report && pdflatex report && pdflatex report
% Numbers final, from the verified empty-cache rerun (output/).
% House style: no long dashes in prose or captions.
\documentclass[10pt,twocolumn]{article}

\usepackage[margin=1.4cm,columnsep=0.6cm]{geometry}
\usepackage{newtxtext,newtxmath} % Times-like text and math
\usepackage{graphicx}
\usepackage{booktabs}
\usepackage{natbib}
\usepackage{enumitem}
\usepackage[hidelinks]{hyperref}

\graphicspath{{../output/plots/}}

% APA-flavoured compact headings.
\usepackage{titlesec}
\titleformat{\section}{\normalsize\bfseries\centering}{}{0pt}{}
\titleformat{\subsection}{\normalsize\bfseries\itshape}{}{0pt}{}
\titlespacing{\section}{0pt}{3pt}{2pt}
\titlespacing{\subsection}{0pt}{4pt}{2pt}

\setlength{\parskip}{0pt}
\setlength{\parindent}{1em}
% Tight float spacing: five column-width figures must fit in two pages.
\setlength{\textfloatsep}{7pt}
\setlength{\floatsep}{6pt}
\setlength{\intextsep}{6pt}
\usepackage[font=footnotesize,labelfont=it]{caption}
\captionsetup{skip=1pt}
\setlist[itemize]{leftmargin=1.0em,itemsep=0pt,parsep=0pt,topsep=1pt}
\setlength{\bibsep}{0pt} % natbib: compact reference list

\title{\large\bfseries Fluent, Generic, and Unfaithful: Evaluating LLM Human
Simulations Against Real Consumer Answers}
\author{\normalsize Dr.\ Chakresh Kr.\ Singh\thanks{Prepared for Roland Berger.}}
\date{}

\begin{document}
\maketitle

\section{The Problem}
A Beauty \& Wellbeing company built LLM-based simulations of individual
consumers from a single open-ended interview, following
\citet{park2024generative}. Can those simulations replace real respondents'
open-ended answers? We treat the evaluation as a \emph{risk control}: where
can simulated answers be trusted for decisions, and where not. Ground truth
is the real humans' answers to a second interview
(3 people $\times$ 10 questions, one human and one AI answer per pair).

\section{Our Outcome}
The simulations are \textbf{not reliably distinguishable in blind reading,
generic rather than individual, and unfaithful} to what the real people
said.

\textbf{Not reliably distinguishable.} Blind judges (one answer, no label,
randomised order) discriminate human from AI at only 0.55 / 0.58 / 0.67
against 0.50 chance. All three judges could better identify human answers than AI (Figure 1)
but could not agree reliably per item. Krippendorff's
$\alpha$ \citep{krippendorff2004content} is chance-corrected agreement
among the three judges on identical items (1 perfect, 0 chance); here
$\alpha = 0.396$, a weak-confidence result.

\textbf{Generic.} When asked whether the AI answer is recognisably \emph{this}
person rather than a generic segment member, the judges could not conclude deterministically.

\textbf{Unfaithful.} These numbers come from the deterministic calculations using encoders. Coverage (how much of what the person said the simulation
kept) is 16.5\%; fidelity (how much of what the simulation said the person
backs) is 4.4\%, with 3 outright contradictions across 30 pairs. The measure is strict by
construction (sentence-level entailment), so absolute values are
understated. 

The combination of the three KPIs makes inference unreliable: detectability judges the \emph{form} of one text alone; coverage and fidelity compare \emph{content} against the
person's real answer. Fluent in-domain text passes blind reading while its
specifics are false, which is why a measurement layer is required.

\textbf{Exploratory judgements agree.} Three judge families each scored the
same 12 held-out items (q7 to q10 for every person), so one verdict is one
judge on one question. Personality capture: 7 of 36 verdicts yes (4, 2, 1
per judge; Figure~\ref{fig:q1}). Content capture: 9 of 36 verdicts yes
(4, 3, 2 per judge; Figure~\ref{fig:q2}). Zero-shot human-or-AI accuracy is
0.58 to 0.78 per judge and asymmetric: every judge recognises humans better
(0.77 to 0.90) than it catches simulations (0.30 to 0.67;
Figure~\ref{fig:q0}).

\textbf{The style gap, measured.} On the deterministic features of
Figure~\ref{fig:style}, the AI writes 55 words to the human's 31, with far
richer vocabulary (MTLD 95 vs.\ 56; \citealp{mccarthy2010mtld}), hedges
less (2.35 vs.\ 5.57 per 100 words), and is strongly positive where humans
are near neutral (VADER mean 0.68 vs.\ 0.28; \citealp{hutto2014vader}).
The simulations sell; the humans shrug.

\begin{figure}[t]
  \centering
  \includegraphics[width=0.52\columnwidth]{eda_q0_accuracy.png}
  \caption{Zero-shot detection accuracy per judge model. Bars show the
  share of the 60 answers each judge labelled correctly, overall and by
  true side; the dashed line marks the 50\% coin-flip level. All three
  judges recognise humans better than they catch simulations.}
  \label{fig:q0}
\end{figure}
\begin{figure}[t]
  \centering
  \includegraphics[width=0.52\columnwidth]{eda_q1_personality.png}
  \caption{Personality-capture verdicts on the held-out questions (q7 to
  q10). Bars count yes verdicts per person, out of 4 per judge model. Only
  7 of 36 verdicts find the person's voice in the AI answer, and most
  concentrate on one person.}
  \label{fig:q1}
\end{figure}
\begin{figure}[t]
  \centering
  \includegraphics[width=0.52\columnwidth]{eda_q2_content.png}
  \caption{Content-capture verdicts on the held-out questions (q7 to q10).
  Bars count yes verdicts per person, out of 4 per judge model. Only 9 of
  36 verdicts find the human answer's content in the AI answer, almost all
  for a single person.}
  \label{fig:q2}
\end{figure}
\begin{figure}[t]
  \centering
  \includegraphics[width=0.52\columnwidth]{eda_style_features.png}
  \caption{Deterministic style features, measured on every answer. Each
  panel compares the human mean (blue) with the AI mean (green) for one of
  the seven features, over 30 answers per side. The AI writes longer, more
  varied, less hedged and far more positive text than the people it
  simulates.}
  \label{fig:style}
\end{figure}
\section{Methodology}
We started from evaluation criteria where each criterion led to one
metric. 

(a) Does it say what the person said: content coverage and fidelity
via NLI. 

(b) Does it sound like a person, and like \emph{this} person: style
features plus an individuation judge. 

(c) Could anyone tell: blind
believability. 

(d) Do the two sides talk about the same themes: a topic
map. 

Below is the detailed methodology with its strengths and weakness.
\begin{itemize}
\item \textbf{Method choice:} sentence-level NLI in both directions with a
DeBERTa cross-encoder \citep{he2021deberta} (coverage and fidelity; SummaC
lineage, \citealp{laban2022summac}). 

\emph{Strength:} deterministic,
auditable, direction-aware; the two directions punish padding and thinness
respectively. 

\emph{Weakness:} strict entailment on conversational text
understates absolute values (domain shift).

\item \textbf{Method choice:} seven deterministic stylometric features per
answer (length, vocabulary, hedging, sentiment; hedging lexicon from
\citealp{hyland2005metadiscourse}). 

\emph{Strength:} fully deterministic,
published lexicons, reproducible. 

\emph{Weakness:} surface-level; cannot
see meaning.

\item \textbf{Method choice:} LLM judge panel of three model families
(gpt-oss:120b, deepseek-v4-flash, glm-5.2) scoring identical items; judges
output only booleans or bounded scores; aggregates computed in code.

\emph{Strength:} family diversity counters self-preference; head-to-head
comparable; agreement measurable. 

\emph{Weakness:} judges are themselves
LLMs; no human raters; the believability construct is weak at
$\alpha = 0.396$.

\item \textbf{Method choice:} blind believability (one unlabelled answer,
randomised order). 

\emph{Strength:} blinding plus randomisation; the
discrimination rate is computed in code. 

\emph{Weakness:} low inter-judge
agreement; judges biased toward calling text human.

\item \textbf{Method choice:} individuation with a persona sketch (the
person's other nine answers). 

\emph{Strength:} the judge sees the person's
own words; our most reliable result ($\alpha = 0.615$). 

\emph{Weakness:}
sketch quality bounded by nine short answers.

\item \textbf{Method choice:} seeded BERTopic topic map over pooled
human/AI sentences. 

\emph{Strength:} targets the population-structure gap
the reference paper concedes. 

\emph{Weakness:} degenerate at this $n$;
reported as not computable.

\item \textbf{Method choice:} fixed explainable split (q7 to q10 held out
for every person), lowercased judge inputs, full disk cache.

\emph{Strength:} explainable, reproducible; orthography confound removed.
\emph{Weakness:} lowercasing also removes any real stylistic case signal.
\end{itemize}

\section{Caveats and Limitations}
\begin{itemize}
\item \textbf{$n=30$ pairs.} We could not test significance
level clearly and had to aggregate at human level. This is because
per-category slices were too small (3 to 15 pairs). 
\item \textbf{Orthography confound, found and corrected.} Run-1 personality
verdicts keyed on lowercase and typos, transcription artefacts rather than
voice; judge inputs were lowercased in subsequent runs
\item \textbf{Topic-level seggregations not computable at this $n$} The
first BERTopic fit collapsed into one stopword mega-topic. We corrected for 
stopwords and used mpnet embeddings with \texttt{n\_neighbors} = 10).  This still led to generic topics with no clear distinction. A similarity-graph alternative was rejected as it would be sensitive to similarity measure thresholds.
\item \textbf{Judge layers are not stable} across cloud reruns despite
temperature 0 and fixed seeds, while deterministic layers reproduce
exactly. 
\item \textbf{No test-retest ceiling.} We could not test humans stability as they were normalised in \citet{park2024generative}. 
\item \textbf{Generation prompts unknown.} Style deviations may be prompt
artefacts (chain-of-thought, expert reflections). These could not be tested
\end{itemize}

\section{Recommendation}
Use the simulations for ideation, hypothesis generation, and
question piloting: tasks that need plausible consumer-shaped text.

Do not use hem for preference measurement, verbatim-style
quotes, or individual-level insight.

The simulations cannot be trusted as substitutes for these consumers'
answers. They are not reliably distinguishable from human answers in blind
reading, which makes their failures invisible to a reader; they revert to
a generic segment member rather than the individual; and almost nothing
they say is backed by what the person actually said, including three
outright contradictions. Their fluency is precisely what makes them
dangerous: nothing warns the reader that the content is wrong. A distinct
but supporting observation is that the LLM components of the evaluation
itself, the judges, were not run-to-run stable despite temperature 0 and
fixed seeds; if LLMs cannot judge consistently, relying on them to
\emph{be} the consumer is a larger leap. The stability of the simulations
themselves was not testable with a single generation per question and is
our first data request.

\bibliographystyle{apalike}
{\footnotesize
\bibliography{references}}

\end{document}
